% !TeX program = xelatex
\documentclass{article}
\usepackage{kotex}
\usepackage[a4paper, left=3cm, right=3cm, top=2.5cm, bottom=2.5cm]{geometry}
\usepackage{amsmath, amssymb}
\usepackage{tcolorbox}
\usepackage{caption}
\usepackage{setspace}
\usepackage{xcolor}
\usepackage{titlesec}
\usepackage{tikz}
\usetikzlibrary{arrows.meta, calc}

\setstretch{1.3}
\renewcommand{\figurename}{그림}
\titleformat{name=\section,numberless}
  {\normalfont\large\bfseries\color{blue}}{}
  {0pt}{}

\title{2.5 쌍대 구면 삼각형}
\date{}
\author{}

\begin{document}
\maketitle

\section*{2.5 쌍대 구면 삼각형}

쌍대(혹은 ``극'') 삼각형이 중요한 이유는 꼭짓점은 원래 삼각형의 변에 대응되고
변은 원래 삼각형의 꼭짓점에 대응되기 때문이다.
따라서 원래 삼각형의 변에 대한 정리는 쌍대 삼각형의 꼭짓점에 대한 정리를 공짜로 제공하고,
원래 삼각형의 꼭짓점에 대한 정리는 쌍대 삼각형의 변에 대한 정리를 제공한다.
이러한 ``쌍대성''을 사용하면 우리의 수고를 반으로 줄일 수 있다.
쌍대성은 수학의 여러 분야에서 사용된다.

\begin{tcolorbox}[
  colback=teal!4, colframe=teal!65!black,
  title=정의 2.5.1 (쌍대 구면 삼각형),
  fonttitle=\bfseries, boxrule=0.5mm, arc=3mm]

$\triangle ABC$를 구면 삼각형이라 하자.
\textbf{쌍대 삼각형} ${}^*\!\triangle ABC$는
$\vec{A}^*,\,\vec{B}^*,\,\vec{C}^*$가 다음과 같은 구면 삼각형 $\triangle A^*B^*C^*$이다.
\[
  \vec{A}^* = \frac{\vec{B}\times\vec{C}}{R\sin\angle a},\qquad
  \vec{B}^* = \frac{\vec{C}\times\vec{A}}{R\sin\angle b},\qquad
  \vec{C}^* = \frac{\vec{A}\times\vec{B}}{R\sin\angle c}
\]
여기서 $R$은 구의 반지름이고, $O$가 구의 중심일 때
$\vec{A}^* = \overrightarrow{OA^*}$, $\vec{B}^* = \overrightarrow{OB^*}$,
$\vec{C}^* = \overrightarrow{OC^*}$이다.

\end{tcolorbox}

위의 식에서 상수 $R\sin\angle a$, $R\sin\angle b$, $R\sin\angle c$는 벡터의 크기를 조절하기 위한 것으로
$|\vec{A}^*| = |\vec{B}^*| = |\vec{C}^*| = R$이 된다.
(두 벡터 $\vec{V}$와 $\vec{W}$에 대해 $|\vec{V}\times\vec{W}| = |\vec{V}||\vec{W}|\sin\angle(\vec{V},\vec{W})$이고,
보조정리 2.3.1에 의해 $\angle a = \angle(\vec{B},\vec{C})$,
$\angle b = \angle(\vec{C},\vec{A})$, $\angle c = \angle(\vec{A},\vec{B})$임을 기억하자.)

$\vec{A}^*$는 원래 삼각형의 변 $a$를 포함하는 평면 $\overline{BCO}$에 수직이다.
비슷하게 $\vec{B}^*$는 $b$를 포함하는 평면에 수직이고 $\vec{C}^*$는 $c$를 포함하는 평면에 수직이다.
이러한 의미로 쌍대 삼각형의 꼭짓점은 원래 삼각형의 변에 대응된다.
다음 정리는 동일한 의미로 쌍대 삼각형의 변이 원래 삼각형의 꼭짓점에 대응됨을 보여준다.
마지막으로 따름정리 2.5.3에 의해 쌍대 삼각형의 각각의 각은 원래 삼각형의 대응되는 각과 보각$^{10)}$ 관계이다.

\begin{tcolorbox}[
  colback=red!4, colframe=red!60!black,
  title=정리 2.5.1,
  fonttitle=\bfseries, boxrule=0.5mm, arc=3mm]

$A^*$, $B^*$, $C^*$를 정의 2.5.1에서의 쌍대 삼각형의 꼭짓점이라 하자.
$a^*$를 $A^*$의 대변, $b^*$를 $B^*$의 대변, $c^*$를 $C^*$의 대변이라 하자.
그러면 다음이 성립한다.
\[
  \vec{A} = s\!\left(\frac{\vec{B}^*\times\vec{C}^*}{R\sin\angle a^*}\right),\qquad
  \vec{B} = s\!\left(\frac{\vec{C}^*\times\vec{A}^*}{R\sin\angle b^*}\right),\qquad
  \vec{C} = s\!\left(\frac{\vec{A}^*\times\vec{B}^*}{R\sin\angle c^*}\right)
\]
단, $s = \pm 1$로서 $\det[\vec{A},\vec{B},\vec{C}]$의 부호이다.$^{11)}$

\end{tcolorbox}

\textbf{증명.}\quad
$\vec{B}^*\times\vec{C}^*$를 계산하면 다음과 같다.
\[
  \vec{B}^*\times\vec{C}^*
  = \left(\frac{\vec{C}\times\vec{A}}{R\sin\angle b}\right)
    \times\left(\frac{\vec{A}\times\vec{B}}{R\sin\angle c}\right)
  = \frac{(\vec{C}\times\vec{A})\times(\vec{A}\times\vec{B})}{R^2\sin\angle b\sin\angle c}
  \tag{2.15}
\]
아래와 같이 주어지는 벡터 항등식을 기억하자.
\[
  \vec{X}\times(\vec{Y}\times\vec{Z}) = (\vec{X}\cdot\vec{Z})\vec{Y} - (\vec{X}\cdot\vec{Y})\vec{Z}
\]
이를 사용하여 식 (2.15)에서 우변의 분자를 계산하면 다음과 같다.
\begin{align*}
  (\vec{C}\times\vec{A})\times(\vec{A}\times\vec{B})
  &= \bigl((\vec{C}\times\vec{A})\cdot\vec{B}\bigr)\vec{A}
   - \bigl((\vec{C}\times\vec{A})\cdot\vec{A}\bigr)\vec{B} \\
  &= \det[\vec{A},\vec{B},\vec{C}]\,\vec{A} + \vec{0}
\end{align*}
그러므로 다음을 얻는다.
\[
  \vec{B}^*\times\vec{C}^*
  = \left(\frac{\det[\vec{A},\vec{B},\vec{C}]}{R^2\sin\angle b\sin\angle c}\right)\vec{A}
\]
이로부터 $\vec{A}$는 다음과 같다.
\[
  \vec{A} = \left(\frac{R^2\sin\angle b\sin\angle c}{\det[\vec{A},\vec{B},\vec{C}]}\right)
  (\vec{B}^*\times\vec{C}^*)
\]
한편, $0°<\angle a^*,\angle b,\angle c < 180°$이므로
$0 < \sin\angle a^*,\sin\angle b,\sin\angle c$이다.
이로부터 다음을 얻는다.

$\vec{A}$는 $s\!\left(\dfrac{\vec{B}^*\times\vec{C}^*}{R\sin\angle a^*}\right)$과 같은 방향을 가리킨다.

또한, $|\vec{B}^*| = |\vec{C}^*| = R$이고 $\angle(\vec{B}^*,\vec{C}^*) = \angle a^*$이므로
\[
  \left|\frac{\vec{B}^*\times\vec{C}^*}{R\sin\angle a^*}\right|
  = \frac{R^2\sin\angle a^*}{R\sin\angle a^*} = R
\]
따라서 $|\vec{A}| = R$이므로 다음이 성립한다.
\[
  \vec{A} = s\!\left(\frac{\vec{B}^*\times\vec{C}^*}{R\sin\angle a^*}\right)
\]
다른 등식도 비슷한 방법에 의해 유도된다. \hfill $\square$

\smallskip
\noindent$^{10)}$\,\small 두 보각은 합하면 $180°$가 된다.\normalsize\\
\noindent$^{11)}$\,\small $\det[\vec{A},\vec{B},\vec{C}]$는 $\vec{A},\vec{B},\vec{C}$를 행벡터로 갖는 $3\times 3$ 행렬의 결정자(행렬식)이다.\normalsize

\begin{tcolorbox}[
  colback=red!4, colframe=red!60!black,
  title=따름정리 2.5.1 (쌍대 삼각형의 쌍대),
  fonttitle=\bfseries, boxrule=0.5mm, arc=3mm]

\[
  {}^*({}^*\!\triangle ABC) =
  \begin{cases}
    \triangle ABC, & \det[\vec{A},\vec{B},\vec{C}] \ge 0 \\[4pt]
    \triangle(-A)(-B)(-C), & \det[\vec{A},\vec{B},\vec{C}] \le 0
  \end{cases}
\]

\end{tcolorbox}

\begin{tcolorbox}[
  colback=red!4, colframe=red!60!black,
  title=따름정리 2.5.2,
  fonttitle=\bfseries, boxrule=0.5mm, arc=3mm]

${}^*({}^*\!\triangle ABC) \cong \triangle ABC$

\end{tcolorbox}

\textbf{증명.}\quad
따름정리 2.5.1에 의해 ${}^*({}^*\!\triangle ABC)$는 $\triangle ABC$와 같거나
구의 중심에 대한 $\triangle ABC$의 반사와 같다. \hfill $\square$

\begin{tcolorbox}[
  colback=red!4, colframe=red!60!black,
  title=따름정리 2.5.3,
  fonttitle=\bfseries, boxrule=0.5mm, arc=3mm]

$A^*$, $B^*$, $C^*$를 정의 2.5.1에서의 쌍대 삼각형의 꼭짓점이라 하자.
$a^*$를 $A^*$의 대변, $b^*$를 $B^*$의 대변, $c^*$를 $C^*$의 대변이라 하자.
그러면 다음이 성립한다.
\begin{align}
  \angle A + \angle a^* = \angle A^* + \angle a &= 180°, \notag\\
  \angle B + \angle b^* = \angle B^* + \angle b &= 180°, \tag{2.16}\\
  \angle C + \angle c^* = \angle C^* + \angle c &= 180° \notag
\end{align}

\end{tcolorbox}

\textbf{증명.}\quad
정의 2.5.1과 보조정리 2.3.1로부터 다음을 얻는다.
\begin{align*}
  \angle a^* &= \angle(\vec{B}^*,\vec{C}^*) \\
             &= \angle(\vec{C}\times\vec{A},\;\vec{A}\times\vec{B}) \\
             &= \angle(-\vec{A}\times\vec{C},\;\vec{A}\times\vec{B}) \\
             &= 180° - \angle(\vec{A}\times\vec{C},\;\vec{A}\times\vec{B}) \\
             &= 180° - \angle A
\end{align*}
따라서 다음이 성립한다.
\[
  \angle A + \angle a^* = 180°
\]
이 결과를 $\triangle A^*B^*C^*$의 쌍대 $\triangle A^{**}B^{**}C^{**}$에 적용하면 다음을 얻는다.
\begin{equation}
  \angle A^* + \angle a^{**} = 180° \tag{2.17}
\end{equation}
여기서 $a^{**}$는 $A^{**}$의 대변이다.
한편, 따름정리 2.5.2에 의해 $\triangle A^{**}B^{**}C^{**}$은 $\triangle ABC$와 합동이므로
$\angle a^{**} \cong \angle a$이다.
그러므로 식 (2.17)에 의해 다음이 성립한다.
\[
  \angle A^* + \angle a = 180°
\]
이로서 식 (2.16)의 첫 번째 식을 증명하였다.
다른 두 식도 비슷한 방법을 적용하면 증명된다. \hfill $\square$

\subsection*{연습문제 2.5}

\textbf{문제 2.5.1}\quad
a) $\mathbb{R}^3$에서 중심이 $(0,0,0)$인 단위 구에 있는 세 점 $A$, $B$, $C$가 다음과 같다.
\[
  A = (0,\,0,\,1),\qquad
  B = \!\left(\frac{\sqrt{2}}{2},\;0,\;\frac{\sqrt{2}}{2}\right),\qquad
  C = \!\left(0,\;\frac{\sqrt{3}}{2},\;\frac{1}{2}\right)
\]
이 때 $\triangle ABC$의 쌍대 삼각형의 꼭짓점을 찾아라.

b) $\triangle ABC$와 쌍대 삼각형에 대해 직접 계산을 통해 정리 2.5.1,
따름정리 2.5.1 그리고 따름정리 2.5.3이 성립함을 확인하여라.

c) 단위 구에 $\triangle ABC$와 ${}^*\!\triangle ABC$를 그려라.

\end{document}
